\subsection{An infeasible strict local minimum in a feasible singleton instance}

The all-inactive stationary point of the unweighted encoding already shows
that stationarity is not a feasibility certificate.  A positive weighting
of the same residuals gives a stronger counterexample: a strict positive
local minimum even when a one-state realization exists.  This example is
for the weighted encoding, not a claim about the local minimality of the
unweighted singleton objective.

Let $J,E,S,I$ and the local hidden carrier all be singletons, take budget
one, and prescribe $A=w=R=1$.  Write the slot coordinates as
$(d,a,b,\rho,\sigma,\tau,\eta)$.  Define
\begin{align*}
 F={}&100\bigl[(d^2-d)^2+(1-d-\eta^2)^2
 +(a-\rho^2)^2+(d-a-\sigma^2)^2\\
 &\hspace{34mm} +(b-\tau^2)^2+2(b-d)^2\bigr]
 +(ad-1)^2+2(ab-1)^2.
\end{align*}
The two copies of $b-d$ are the continuation normalization and the
singleton local-marginal equation.  The three data residuals impose
initialization, prior, and output, respectively.  Multiplying each
structural residual by $10$ before squaring leaves the exact zero set
unchanged; the objective remains a quartic polynomial sum of squares.

At
\[
 P_0=(0,0,0,0,0,0,1),\qquad
 P_1=(1,1,1,1,0,1,0),
\]
one has $F(P_0)=3$, $\nabla F(P_0)=0$, and $F(P_1)=0$.
In fact, $P_0$ is a strict local minimum.  To prove this, set
$L=a-\rho^2$, $G=d-a-\sigma^2$, and $N=b-d$.
For $|d|<1/2$,
\[
 (d^2-d)^2\geq d^2/4,\qquad
 a^2\leq 2d^2+2G^2+L^2,\qquad
 b^2\leq 2d^2+2N^2.
\]
For the middle inequality, if $a<0$ then $a^2\leq L^2$;
if $a\geq0$, the equation $a+\sigma^2=d-G$ gives
$a\leq |d|+|G|$.  The data contribution satisfies
\begin{align*}
 (ad-1)^2+2(ab-1)^2-3
 &=a^2d^2+2a^2b^2-2ad-4ab\\
 &\geq a^2d^2+2a^2b^2-d^2-3a^2-2b^2.
\end{align*}
Combining the inequalities yields the explicit lower bound
\begin{align*}
 F-3\geq{}&14d^2+94G^2+97L^2+96N^2+a^2d^2+2a^2b^2\\
 &+100\bigl[(1-d-\eta^2)^2+(b-\tau^2)^2+N^2\bigr].
\end{align*}
Every term on the right is nonnegative.  Equality forces
$d=L=G=N=0$, then $a=b=\rho=\sigma=0$, then $\tau=0$
and $\eta^2=1$.  In the open neighborhood $|d|<1/2$, $\eta>0$,
equality therefore occurs only at $P_0$.
Thus a local optimizer can stop at a strict positive minimum although an
exact feasible solution exists.  The proof uses the quartic terms, not
just the Hessian: the full Hessian has flat $\rho,\sigma,\tau$
directions at $P_0$.

\paragraph{Independent exact checks.}
Symbolic differentiation gives all seven partial derivatives zero at
$P_0$.  The Hessian principal block in $(d,a,b,\eta)$ is
\[
 \begin{pmatrix}
 1000&-202&-400&400\\
 -202&400&-4&0\\
 -400&-4&600&0\\
 400&0&0&800
 \end{pmatrix}.
\]
Its leading principal minors are
$1000$, $359196$, $150855200$, and $82286720000$.
These provide regression checks on the specified coefficients, while the
inequality above establishes strict local minimality.

\subsection{Small residuals and attainment must be distinguished}

The coercive polynomial encoding has an attained global minimum for every
fixed input instance.  In particular, for that encoding, infimum zero is
equivalent to existence of an exact zero.  This excludes an escape to
infinity, but it does not turn a small computed value or a local minimum
into an exact feasibility decision.

For a quantitative example, take $0<\varepsilon<1$,
$J=\{*\}$, $H_1=H_2=\{0,1\}$,
$A_1=A_2=(1/2,1/2)$, and
$E=\{(0,0),(1,1)\}$.  Let $S=\{0,1\}$ be a shared coarse
carrier, with both local projections the identity, and set
\[
 B_1=\begin{pmatrix}1&0\\0&1\end{pmatrix},\qquad
 B_2=\begin{pmatrix}1-\varepsilon&\varepsilon\\
                   \varepsilon&1-\varepsilon\end{pmatrix}.
\]
Set $Q$ to the identity and $w=R=(1/2,1/2)$.
All aggregate observed distributions agree.  Nevertheless, each allowed
$K_h$ is empty because a shared continuation row would have to equal two
different prescribed vectors.  No realization exists at any budget.

Select two slots, one for each allowed tag, with root weights $1/2$ and
continuation rows
\[
 b_0=(1-\varepsilon/2,\varepsilon/2),\qquad
 b_1=(\varepsilon/2,1-\varepsilon/2).
\]
All constraints except local row preservation hold exactly.  The total
squared local-row residual is $2\varepsilon^2$.  Thus infeasible instances
can have arbitrarily small residual values.  This is a family of changing
inputs, not a fixed instance with unattained infimum zero.  For each fixed
positive $\varepsilon$, coercivity and infeasibility imply a positive
attained global minimum.  There is no universal positive tolerance that
certifies exact feasibility for every input in this family.

\paragraph{Traps excluded by the chosen parameterization.}
A strictly positive parameterization would miss valid boundary solutions.
For example,
\[
 \beta(t)=\left(\frac{e^t}{1+e^t},\frac{1}{1+e^t}\right)
\]
has squared error $2/(1+e^t)^2$ relative to $(1,0)$, with infimum zero but
no finite parameter attaining it.  Square variables retain the boundary.
Separately, omitting the normalization and gate bounds can permit
polynomial escape: $(ab-1)^2+a^2$ on $a,b\geq0$ has no zero but equals
$1/n^2$ at $(a,b)=(1/n,n)$.  These are counterexamples to those alternative
formulations, not to the coercive activity-slot encoding.
